with a @@M0@@-ring structure, relations (2.4.2) passing to the quotient.

Every intersection of @@M1@@-ideals is a @@M2@@-ideal. In particular, there exists a smallest @@M3@@-ideal containing a given family of elements @@M4@@; by relations (2.4.2), it is the (ordinary) ideal generated by the @@M5@@, @@M6@@.

Again let @@M7@@ be a @@M8@@-ring, @@M9@@ a family of indeterminates, @@M10@@, and let @@M11@@ be a family of elements of @@M12@@: let @@M13@@ be the @@M14@@-ideal generated by the @@M15@@.

\sgaStableTarget{sga6.e05.def.4-6}\textbf{Definition 4.6.} \uline{The} @@M16@@-\uline{ring} @@M17@@ \uline{will be called the} @@M18@@-ring generated over @@M19@@ by the family of generators @@M20@@ subject to the relations @@M21@@.

We can also specify relations between the @@M22@@, @@M23@@ and make the same construction, observing that an ideal is a @@M24@@-ideal if and only if it is stable under the operations @@M25@@, @@M26@@, as follows from relations 3.2.

Let @@M27@@ be a @@M28@@-@@M29@@-algebra, @@M30@@ a family of elements of @@M31@@. Since the @@M32@@ are polynomials in the indeterminates @@M33@@, @@M34@@, we can substitute the @@M35@@, @@M36@@ for these indeterminates; if the resulting element is zero, we say that the @@M37@@ satisfy the relations @@M38@@. The @@M39@@-@@M40@@-algebra @@M41@@ is then characterized by the following universal property:

\sgaStableTarget{sga6.e05.prop.4-7}\textbf{Proposition 4.7.} \uline{For every} @@M42@@-@@M43@@-\uline{algebra} @@M44@@ \uline{and every family} @@M45@@ \uline{of elements of} @@M46@@, \uline{satisfying the relations} @@M47@@, \uline{there exists exactly one} @@M48@@-@@M49@@-\uline{homomorphism} @@M50@@ \uline{such that for every} @@M51@@, we have @@M52@@.

Indeed, the kernel of the homomorphism @@M53@@ defined in 4.2 contains the @@M54@@, hence @@M55@@ since it is a @@M56@@-ideal; @@M57@@ therefore passes to the quotient and defines
